% Eje 2 (ejemplo: Álgebra y funciones)

    \item En el siguiente gráfico se representa una función afín $f$.
        \begin{center}
            \begin{tikzpicture}[scale=0.55]
                \draw[relleno] (-1, -1) grid (5, 9);
                \draw[-latex] (-1, 0) -- (5.5, 0) node[right] {$X$};
                \draw[-latex] (0, -1) -- (0, 9.5) node[left] {$Y$};
                \draw[very thick, acento] (-0.5, 1) -- (3.5, 9);
                \fill (0, 2) circle (3pt) node[left] {$2$};
                \fill (3, 8) circle (3pt);
                \draw[dashed] (0, 8) node[left] {$8$} -- (3, 8) -- (3, 0) node[below] {$3$};
            \end{tikzpicture}
        \end{center}

    ¿Cuál de las siguientes expresiones corresponde a $f(x)$?
    \begin{enumerate}
        \item $2x+2$ % <- correcta
        \item $x+2$
        \item $2x+8$
        \item $3x+2$
    \end{enumerate}
